\chapter{求解方法}\label{sec:solution_methodology}

\section{总体求解算法}
\label{sec:algorithm}

如第~\ref{sec:complexity}~节所讨论的，完整的\problemName 公式是NP-难的，由于时空网络和行程选择层的组合增长，在实践中变得计算困难。虽然标准的动态离散化发现（DDD）方法高效地压缩物理网络节点和弧，但它缺乏处理高度耦合的基于行程的旅客选择行为的机制。为弥合这一差距，我们提出行程聚合与动态细化（\algoNotation）。与标准DDD的根本区别在于，\algoNotation 将动态聚合从物理网络维度扩展到逻辑旅客维度。通过引入路径签名的概念，\algoNotation 将数百万条组成行程分组为紧凑的超级行程集合，而不丧失选择模型的结构完整性和解的最优性。通过这种双重聚合建立紧的下界、恢复逻辑有效的上界，以及精确插入缺失的时间点以纠正物理时序违规和行程级虚假收入，该算法框架保证收敛到精确的全局最优解。

本节概述了这一完整的精确求解框架，为三个基本算法构建模块奠定基础：提供有效下界的松弛模型（第~\ref{sec:lbm}~节）、作为修复算法的上界生成过程（第~\ref{sec:ub}~节）以及确保在迭代中收敛到全局最优的离散化细化策略（第~\ref{sec:refinement}~节）。

\subsection{初始离散化}\label{sec:initial_discretization}

\algoNotation 算法从一个刻意粗糙的初始离散化$\discretization$开始，以保持松弛模型的小规模。然而，初始步长的选择对松弛质量有深远影响。如第~\ref{sec:lbm}~节将详细介绍的，用于在松弛网络中构造服务弧的下取整映射规则（性质~\ref{def:relax.flightArc}）在网格相对于航段旅行时间过于粗糙时可能产生物理上无意义的逆弧。为从一开始就消除这种病态，我们采用一种初始离散化，其中每个节点$i \in \airportSet$获得一个根据其最短入站旅行时间校准的个体步长$\initStep_i$。精确的构造及其理论证明在第~\ref{sec:init_disc_detail}~节中给出。

\subsection{迭代框架}
我们从初始离散化开始，迭代求解下界模型以获得松弛解及其对应的下界。然后，我们使用松弛解通过基于MIP的修复算法生成上界。如果上界和下界之间的间隙在预定义阈值$\epsilon>0$内，则终止算法。否则，我们根据当前解细化离散化并重复该过程，这确保收敛到完整模型的全局最优解。算法总结在算法~\ref{alg:refinement}中。

\begin{algorithm}[ht]
\caption{Iterative Refinement Algorithm for $\problemNotation$}
\label{alg:refinement}
\begin{algorithmic}[1]
\State $LB \gets -\infty$.
\State $UB \gets +\infty$.
\State $\solution^{*} \gets \text{null}$.
\ForAll{$i\in\airportSet$}
\Comment{Initialize (Section~\ref{sec:initial_discretization})}
    \State $\initStep_i\gets\min_{j:(j,i)\in\segSet}\{\travelTime_{ji}\}$
    \State $\discretization_i\gets\{t\mid t=\planningStart+k\,\initStep_i,\, k \in \mathbb{Z}_{\ge0}, t < \planningEnd\}$
\EndFor
\State $iter\gets 0$
\Loop
    \State $iter \gets iter + 1$
    \State $\solution, LB \gets\Call{ATF}{\arcSet}$
    \Comment{Update Lower Bound}
    \State $\solution^{*}_{new}, UB_{new} \gets \Call{UB\_Model}{\solution}$
    \If{$UB_{new} < UB$}
        \State $\solution^{*}, UB \gets \solution^{*}_{new}, UB_{new}$
        \Comment{Update Upper Bound}
    \EndIf
    \If{$UB - LB \le \epsilon \;\textbf{or}\; iter \ge \text{maxIter}$}
        \State \textbf{break}
    \EndIf
    % \ForAll{$i \in \airportSet, t \in \discretization_i, f \in \fleetSet$}
    %     \Comment{Refinement}
    %     \State $Q_{imp} \gets D(i, t, f) - H_{in}(i, t, f) - A_{legal}(i, t, f)$
    %     \If{$Q_{imp} > 0$}
    %         \State $\mathcal{K} \gets A_{illegal}(i, t, f)$
    %         \State Sort $\mathcal{K}$ by flow volume $\fleetAssignVar_{k}$ in descending order
    %         \State $covered \gets 0$
    %         \ForAll{$k \in \mathcal{K}$}
    %             \If{$covered < Q_{imp}$}
    %                 \State $\discretization_i \gets \discretization_i \cup \{t_{\text{dep}}(k) + \travelTime_k\}$
    %                 \State $covered \gets covered + \fleetAssignVar_{k}$
    %             \Else
    %                 \State \textbf{break}
    %             \EndIf
    %         \EndFor
    %     \EndIf
    % \EndFor
    \State $\discretization \gets \Call{Refinement}{\solution, \discretization}$
    \Comment{Refine Discretization}
\EndLoop
\State \Return $\solution, \solution^{*}, LB, UB$
\end{algorithmic}
\end{algorithm}

\section{吸引力平衡预处理}\label{sec:preconditioning_solution}

为稳定优化过程，我们提出一种市场特定的数值预处理技术，称为吸引力平衡。对于每个市场$m \in \marketSet$和旅客类型$p \in \psgTypeSet$，我们在所有吸引力相关参数中识别最小非零绝对吸引力值：
\begin{equation*}
    \attrMin_{mp} = \min \left\{ \left|\outAttr_{mp}\right|, \, \left|\attr_{rp\ell}\right|, \, \left|\adjAttr_{rp\ell}\right| \mid r \in \itinInMarket_m, \, \ell \in \fareClassSet, \, \text{value} \neq 0 \right\}
    \label{eq:attr_min}
\end{equation*}

然后我们应用缩放因子$\attrScale_{mp} = 1.0 / \attrMin_{mp}$来均匀缩放该市场-旅客组合内的所有吸引力参数。由于根据GAM定义$\adjAttr_{rp\ell} = \attr_{rp\ell} - \shadowAttr_{rp\ell}$，我们缩放底层的名义吸引力和影子吸引力：
\begin{equation*}
\begin{aligned}
    \outAttr_{mp} &\leftarrow \attrScale_{mp} \cdot \outAttr_{mp}\\
    \attr_{rp\ell} &\leftarrow \attrScale_{mp} \cdot \attr_{rp\ell} \quad \forall r \in \itinInMarket_m, \, \ell \in \fareClassSet\\
    \shadowAttr_{rp\ell} &\leftarrow \attrScale_{mp} \cdot \shadowAttr_{rp\ell} \quad \forall r \in \itinInMarket_m, \, \ell \in \fareClassSet
\end{aligned}
\end{equation*}
因此，调整吸引力也按相同因子缩放，即$\adjAttr_{rp\ell} \leftarrow \attrScale_{mp} \cdot \adjAttr_{rp\ell}$和$\adjOutAttr_{mp} \leftarrow \attrScale_{mp} \cdot \adjOutAttr_{mp}$。

从数学上看，由于选择约束\eqref{constr:main.choice.ratio}和\eqref{constr:main.choice.limit}在每个市场-旅客组合内对吸引力参数是齐次的，将所有吸引力值乘以正常数$\attrScale_{mp}$严格保持可行域和最优解。为验证这一点，观察约束\eqref{constr:main.choice.ratio}变为：
\begin{equation*}
    (\attrScale_{mp} \cdot \outAttr_{mp})\, \shareVar_{rp\ell} \;\le\; (\attrScale_{mp} \cdot \attr_{rp\ell})\, \outsideShareVar_{mp}
\end{equation*}
这在从两边消去公共正因子$\attrScale_{mp}$后等价于原始约束。类似地，约束\eqref{constr:main.choice.limit}仍然有效，因为系数中的分子和分母按相同因子缩放，保持了比率$\adjAttr_{rp\ell}/\attr_{rp\ell}$和$\adjOutAttr_{mp}/\outAttr_{mp}$。

在计算上，这种预处理确保选择约束中最小的非零系数恰好为$1.0$，有效消除了矩阵病态。这一简单但关键的变换在我们求解框架的下界模型和上界模型组件中都起到催化剂作用，大幅加速了下界的收敛并增强了整体算法稳定性。本文报告的所有实验结果均在所有实例上统一采用了这一吸引力平衡预处理步骤。

一种解决数值不稳定性的替代方法是将$\shareVar_{rp\ell}$替换为$\demand_{mp} \shareVar_{rp\ell}$，其中唯一的$m$满足$r\in \itinInMarket_m$，并从目标函数和约束中移除$\demand_{mp}$。同时，$\shareVar_{rp\ell}$和$\outsideShareVar_{mp}$的域变为$[0, \demand_{mp}]$（或简单地$[0,+\infty]$）。这种方法可以有效消除目标函数和选择约束中的大系数，特别是$\demand_{mp}\fare_{r\ell}$项。新$\shareVar$变量的值表示选择行程$r$在票价$\ell$下的实际旅客数量，而非需求份额，后者通常是$[10^{-5}, 10^{-4}]$中的小分数。

\section{Relaxation Model $\model{\discretization, \itinSetSuper}$ via Discritization and Itinerary Aggregation}
\label{sec:lbm}

To address the tractability issues of the full time-space network, we introduce a time relaxation approach. The central idea of this step is to formulate a relaxed model on a much coarser, partially discretized time-space network, thereby drastically reducing the problem size while still guaranteeing a valid theoretical lower bound for the exact algorithm. This base model ignores certain strict timing feasibility requirements which are subsequently resolved in later algorithm stages. The model proposed in this section is called the Lower Bound Model (LBM) or the relaxation model in the rest of the paper.

% \section*{Understanding $\discretization$}
We consider $\discretization$ as a transformation of $\discretizationFull$. The transformation can be seen as aggregations of time-arcs. 
Every time-arc in $\full{\net}$ have a corresponding arc in $\net$.That means, we have some time-nodes $(i,t)$ in $\discretization$. All time-arcs $\full{\net}$ will exist in $\net$, and their start time, end time, and travel time may change because they can only use time-nodes in $\discretization$. In order to ensure the optimality, we can let the travel time be shorter than what they are. And there may even be inverse arcs whose start time is later than finish time.
The time-expanded network $\net = (\nodeSet, \arcSet)$, where $\nodeSet = \{(i,t)\mid \forall i\in\airportSet, t\in\discretization_i\},\;\arcSet = \holdingArc \cup \flightArc$, constructed from time discretization $\discretization$ satisfies the following properties:
% \[
%     \net = (\nodeSet, \arcSet)
% \]
% \[
%     \nodeSet = \{(i,t)\mid \forall i\in\airportSet, t\in\discretization_i\},\;\arcSet = \holdingArc \cup \flightArc
% \]
% \[
%     \net = \Phi(\netFlat, \discretization) = \Phi(\airportSet,\segSet,\discretization)
% \]
\begin{property}[Domain of Discretization $\discretization$]
    $\forall i\in\airportSet,\,\{\planningStart\}\subseteq\discretization_i\subseteq\discretizationFull_i$
\end{property}
\begin{property}[Nodes $\nodeSet = \Phi_{\nodeSet}(\netFlat, \discretization)$]
    Node set $N$ contains node $(i,\planningStart)\;\forall i\in\airportSet$
\end{property}
\begin{property}[Holding Arcs]
    \label{def:relax.holdingArc}
    Each ground arc $a\in\holdingArc$ is defined as $a=((i,t),(i,t^{+}))$, where
    $$t^{+}=\begin{cases}
            \min\{\bar{t}\mid\bar{t}>t,\bar{t}\in\discretization_i\} & ,\exists\bar{t}>t,\bar{t}\in\discretization_i\\
            \planningStart & ,\text{otherwise}
    \end{cases}$$
\end{property}
\begin{property}[Service Arcs]
    \label{def:relax.flightArc}
    Each flight arc $a\in\flightArc$ is defined as $a=((i,t),(j,t^\prime)),\;\forall (i,j)\in\segSet,t\in\discretization_i$, where
    $$t^\prime=\max\{\bar{t}\mid\bar{t}\leq(t+\travelTime_{ij})\;\text{mod}\;(\planningEnd-\planningStart)+\planningStart,\bar{t}\in\discretization_j\}$$
\end{property}

Denote the arc mapping induced by the discretization transformation as
$\Phi_{\arcSet}:\full{\arcSet}\to \widetilde{\arcSet}$,
where $\widetilde{\arcSet}:=\arcSet\cup\{\bot\}$ augments the partial network with a virtual symbol $\bot$ capturing collapsed holding arcs.
For $a\in\arcSet$ define its pre-image
$\Phi_{\arcSet}^{-1}(a)=\{a'\in\full{\arcSet}\mid\Phi_{\arcSet}(a')=a\}$, and let
$\Phi_{\arcSet}^{-1}(\bot)$ collect the arcs mapped to the virtual element.
Define analogously $\Phi_{\nodeSet}:\full{\nodeSet}\twoheadrightarrow\nodeSet$
and $\Phi_{\nodeSet}^{-1}(i,t)=\{(i,\tau)\in\full{\nodeSet}\mid\Phi_{\nodeSet}(i,\tau)=(i,t)\}$.

For every node $i\in\airportSet$, Property~1 gives $\{\planningStart\}\subseteq\discretization_i\subseteq\discretizationFull_i$,
so we may introduce the projection
\[
    \phi_i:\discretizationFull_i\to\discretization_i,\qquad 
    \phi_i(\tau)=\max\{\bar{t}\in\discretization_i\mid \bar{t}\le \tau\}.
\]
Property~2 yields $\Phi_{\nodeSet}((i,\tau))=(i,\phi_i(\tau))$.
Using Properties~3 and 4, the arc mapping of $(i,\tau)\rightarrow(j,\tau')$ is specified as
\begin{equation}
    \Phi_{\arcSet}(((i,\tau)\rightarrow(j,\tau'))) =
    \begin{cases}
        ((i,\phi_i(\tau))\rightarrow(j,\phi_j(\tau'))), & \text{if } i\neq j,\\[2pt]
        ((i,\phi_i(\tau))\rightarrow(i,\phi_i(\tau'))), & \text{if } i=j,\ \phi_i(\tau)<\phi_i(\tau'),\\[2pt]
        \bot, & \text{if } i=j,\ \phi_i(\tau)=\phi_i(\tau').
    \end{cases}
\end{equation}
Each service arc keeps its tail on the rounded-down partial times and sets its head to the latest partial time point that the travel time is not increased, whereas a holding arc bridges to the next partial node if one exists, and otherwise collapses into $\bot$ when both endpoints project to the same partial node.

\section*{Variables}
\subsection*{Fleet Assignment Variable $\fleetAssignVar_{af} = \fleetAssignVar_{a^{\discretization},f}^{\net}$}
\[
    \fleetAssignVar_{a^{\discretization},f}^{\net} = \sum_{a^{\discretizationFull}:\, \Phi(a^{\discretizationFull}) = a^{\discretization}}\fleetAssignVar_{a^{\discretizationFull},f}^{\full{\net}}
\]
\subparagraph*{Domain}
\begin{equation}
    \fleetAssignVar_{af}\in\mathbb{Z}_{\ge0} \quad \forall a\in \flightArc, \; \forall f\in \fleetSet
    % temp-command \tag{2j}
    \label{constr:relax.domain.x.flight}
\end{equation}
which ensures optimality.
Therefore domains of some $\fleetAssignVar_{a,f}$ need to change.

\section*{Constraints Transformation}
Below we show systematically about the algorithmical influence of applying the partial discretization.
Considering each constraint. (\ref{constr:main.flow}, \ref{constr:main.seg_freq}) keeps. (\ref{constr:main.choice.ratio}, \ref{constr:main.choice.limit}, \ref{constr:main.choice.capacity}, \ref{constr:main.choice.covering}) keeps. The only differences of constraints above are the sets of arcs and nodes are changed to the ones in the partial network. However, some constraints need to be modified to ensure the optimality of the relaxation model.

% (\ref{constr:main.fleet_available}) may have problems when there are inverse arcs. Right hand sides of (\ref{constr:main.cooldown}) need to be the value of the maximum flights of next time periods.

An resource unit may be counted twice when traveling through blue paths. The blue service arcs are not counted directly, but their two blue holding arcs are, each adding 1 to the total.
\begin{center}
    \begin{tikzpicture}[x=1cm,y=0.75cm]
        % axes (time lines) for i and j
        \draw[timeaxis] (0,0) -- (10,0);
        \draw[timeaxis] (0,1) -- (10,1);
        % labels
        \node[left] at (0,0) {$i$};
        \node[left] at (0,1) {$j$};

        % time points
        \foreach \x in {0,1,2,3,4,5,6,7,8,9}{
          \node[dot] (i\x) at (\x+0.3,0) {};
          \node[dot] (j\x) at (\x+0.3,1) {};
        }

        % repeated flight arcs i->j
        \foreach \k in {0,2,3,5,6,7}{
          \draw[flight] (i\k) -- ($(j\k)+(2,0)$);
        }
        \draw[timec,gray!70] (5.8,-0.3) -- (5.8,1.3);

        \draw[timecActive,red!70] (4.3,0) -- (5.3,0);
        \draw[timecActive,red!70] (5.3,0) -- (7.3,1);
        \draw[timecActive,red!70] (7.3,1) -- (8.3,1);
        
        \draw[timecActive,blue!70] (5.3,0) -- (6.3,0);
        \draw[timecActive,blue!70] (6.3,0) -- (8.3,1);
        \draw[timecActive,blue!70] (8.3,1) -- (9.3,1);

        % forward time arrows (just to suggest direction)
        \draw[flight] (10.0,0) -- (10.25,0);
        \draw[flight] (10.0,1) -- (10.25,1);
    \end{tikzpicture}
    \begin{tikzpicture}[x=1cm,y=0.75cm]
        % axes (time lines) for i and j
        \draw[timeaxis] (0,0) -- (10,0);
        \draw[timeaxis] (0,1) -- (10,1);
        % labels
        \node[left] at (0,0) {$i$};
        \node[left] at (0,1) {$j$};

        % time points
        \foreach \x in {0,1,2,3,4,5,6,7,8,9}{
          \node[dot] (i\x) at (\x+0.3,0) {};
        }
        \foreach \x in {0,1,2,3,4,5,9}{
          \node[dot] (j\x) at (\x+0.3,1) {};
        }

        % repeated flight arcs i->j
        \foreach \k in {0,2,3,7}{
          \draw[flight] (i\k) -- ($(i\k)+(2,1)$);
        }
        \foreach \k in {5,6}{
          \draw[flight] (i\k) -- ($(j5)$);
        }

        %
        \draw[timec,gray!70] (5.8,-0.3) -- (5.8,1.3);

        \draw[timecActive,red!70] (4.3,0) -- (5.3,0);
        \draw[timecActive,red!70] (5.3,0) -- (5.3,1);
        \draw[timecActive,red!70] (5.3,1) -- (9.3,1);
        
        \draw[timecActive,blue!70] (5.3,0) -- (6.3,0);
        \draw[timecActive,blue!70] (6.3,0) -- (5.3,1);
        \draw[timecActive,blue!70] (5.3,1.05) -- (9.3,1.05);

        % forward time arrows (just to suggest direction)
        \draw[flight] (10.0,0) -- (10.25,0);
        \draw[flight] (10.0,1) -- (10.25,1);
    \end{tikzpicture}
\end{center}
Let set $\invertArcTc{t_c}$ denotes inverse(except wrapped) arcs across $t_c$:
\begin{equation}
\begin{aligned}
    \invertArcTc{t_c} = \{a=((i,t)\rightarrow(j,t^\prime))\;\forall a\in\flightArc\mid t^\prime\leq t_c<t<t+\travelTime_{ij}<\planningEnd\}
    % temp-command \tag{def$_{\invertArcTc{t_c}}$}
    \label{def:relax.invertArcTc}
\end{aligned}
\end{equation}
Therefore, the fleet availability constraints are modified to
\begin{equation}
    \sum_{a\in \holdingArcTc{\checkTime}} \fleetAssignVar_{af} + \sum_{a\in \flightArcTc{\checkTime}} \fleetAssignVar_{af} - \sum_{a\in \invertArcTc{\checkTime}} \fleetAssignVar_{af} \le \fleetAvail_{f}
    \quad \forall f\in \fleetSet
    % temp-command \tag{2c}
    \label{constr:relax.fleet_available}
\end{equation}

Aggregation may cause numbers of services which are originally far away from each other depart at same or relative closer time in the relaxation model. In order to ensure the optimality of the relaxation model, we need to enable departure-pairs which are legal in the full network.
Right hand sides of (\ref{constr:main.cooldown}) need to be the value of the maximum services of next time periods.
% \subparagraph*{Set of arcs for segment $(i,j)$ at time $t$: $\arcDuringSepTime_{ijt}^\discretization$}
% % \begin{align}
% %     \arcDuringSepTime_{ijt}^\discretization = \{&a=((i,t^\prime)\rightarrow(j,t^{\prime\prime}))\;\forall a\in\flightArc\mid
% %     && t\leq t^\prime<t+\minInterdepTime_{ij}<\planningEnd \cup t\leq t^\prime<\planningEnd<t+\minInterdepTime_{ij}\notag \\
% %     &&& \cup t^\prime<t+\minInterdepTime_{ij}-(\planningEnd-\planningStart)<\planningEnd<t+\minInterdepTime_{ij}\}\notag \\
% %     = \{&a=((i,t^\prime)\rightarrow(j,t^{\prime\prime}))\;\forall a\in\flightArc\mid
% %     && t\leq t^\prime<t+\minInterdepTime_{ij}\notag \\
% %     &&& \cup t^\prime<t+\minInterdepTime_{ij}-(\planningEnd-\planningStart)<\planningEnd<t+\minInterdepTime_{ij}\}
% %     % temp-command \tag{def$_{\arcDuringSepTime_{ijt}^\discretization}$}
% %     \label{def:relax.arcDuringSepTime}
% % \end{align}
% \begin{equation}
%     \arcDuringSepTime_{ijt}^\discretization = \{a=((i,t^\prime)\rightarrow(j,t^{\prime\prime}))\;\forall a\in\flightArc\mid t\leq t^\prime<t+\minInterdepTime_{ij} \}
%     \label{def:relax.arcDuringSepTime}
% \end{equation}
The upper bound of the number of services which can depart in the interval $[t, t+\minInterdepTime_{ij})$, $\departLimit_{ijt}$ is defined as:
\begin{equation}
\begin{aligned}
    \departLimit_{ijt} & = \begin{cases}
        \lceil \frac{t^{+}-t}{\minInterdepTime_{ij}}\rceil & t^{+}-t>0\notag\\
        \lceil \frac{\planningEnd-\planningStart+t^{+}-t}{\minInterdepTime_{ij}}\rceil & t^{+}-t<0
    \end{cases} \\
    & = \lceil \frac{t^{+}-t+\frac{\planningEnd-\planningStart}{2}(\frac{t^{+}-t-|t^{+}-t|}{|t^{+}-t|})}{\minInterdepTime_{ij}}\rceil
    % temp-command \tag{def$_{\departLimit_{ijt}}$}
    \label{def:relax.departLimit}
\end{aligned}
\end{equation}
where
$
    t^{+}=\min\{\bar{t}\in\discretization_i\mid\bar{t}>\max\{\bar{\bar{t}}\mid(i,\bar{\bar{t}},j,\bar{\bar{t}}^\prime)\in\arcDuringSepTime_{ijt}^\discretization\}\}
$
Therefore, the frequency constraints are modified to
\begin{equation}
    \sum_{f\in \fleetSet}\sum_{a\in \arcDuringSepTime_{ijt}} \fleetAssignVar_{af} \le \departLimit_{ijt}
    \quad \forall (i,j)\in \segSet,\; \forall t \in \discretization_i
    % temp-command \tag{2d}
    \label{constr:relax.cooldown}
\end{equation}

\subsection{Topology-Aware Heterogeneous Initial Discretization}\label{sec:init_disc_detail}

The arc-mapping rules defined in Properties~\ref{def:relax.holdingArc}--\ref{def:relax.flightArc} reveal a structural vulnerability of uniform coarse discretizations. Under Property~\ref{def:relax.flightArc}, the arrival time of a service arc on segment $(j,i)$ departing at $t \in \discretization_j$ is set to the latest grid point in $\discretization_i$ no later than $t + \travelTime_{ji}$ (floor mapping). If a uniform step size $\discretizationStep$ is applied across all nodes and $\discretizationStep > \travelTime_{ji}$ for some segment $(j,i)$, the mapped arrival time $t'$ may satisfy $t' \le t$, producing an inverse arc in which a resource unit appears to arrive no later than it departs. As shown in Equation~\eqref{def:relax.invertArcTc} and Constraint~\eqref{constr:relax.fleet_available}, inverse arcs require dedicated correction terms in the fleet availability constraints. More critically, they open up a large region of physically infeasible solutions in the relaxation, severely degrading the lower bound quality.

To eliminate inverse arcs at their source, we adopt a topology-aware heterogeneous initial discretization. Instead of a single global step, each node $i \in \airportSet$ receives an individual initial step size $\initStep_i$ defined as the minimum inbound travel time:
\begin{equation}\label{eq:init_step}
    \initStep_i = \min_{j:\,(j,i)\in\segSet} \left\{ \travelTime_{ji} \right\}.
\end{equation}
For nodes with no inbound segments, we set $\initStep_i = \planningEnd - \planningStart$ as a safe default. The initial time point set for node $i$ is then generated as the arithmetic progression
\begin{equation}\label{eq:init_discretization}
    \discretization_i = \left\{ t \;\middle|\; t = \planningStart + k \cdot \initStep_i,\quad k \in \mathbb{Z}_{\ge 0},\quad t < \planningEnd \right\}.
\end{equation}

\begin{proposition}[Elimination of Inverse Arcs]\label{prop:no_inverse}
    Under the initial discretization defined by \eqref{eq:init_step}--\eqref{eq:init_discretization}, the relaxed network $\net$ contains no inverse arcs, i.e., for every flight arc $a = ((j,t) \to (i,t')) \in \flightArc$ with $(j,i) \in \segSet$, we have $t' > t$.
\end{proposition}
\begin{proof}
    By Property~\ref{def:relax.flightArc}, $t' = \max\{\bar{t} \in \discretization_i \mid \bar{t} \le t + \travelTime_{ji}\}$. Since $\discretization_i$ is an arithmetic progression with common difference $\initStep_i \le \travelTime_{ji}$ (by \eqref{eq:init_step}), the interval $(t,\, t + \travelTime_{ji}]$ has length $\travelTime_{ji} \ge \initStep_i$ and therefore contains at least one grid point of $\discretization_i$. Hence $t' \ge t + \initStep_i > t$ (noting $\initStep_i > 0$ by construction).
\end{proof}

Beyond ensuring physical feasibility, this strategy provides two additional benefits.

\paragraph{Topology-adaptive variable compression.}
The heterogeneous step sizes naturally adapt to the hub-and-spoke structure of real-world service networks. Hub nodes, which receive many short-haul services, obtain a small $\initStep_i$ and hence a finer grid that preserves temporal precision. Spoke nodes served predominantly by long-haul segments receive a large $\initStep_i$ and a correspondingly coarser grid. This self-adaptive behavior keeps the total number of initial time points---and consequently the number of decision variables---compact without sacrificing the physical validity of the network, avoiding out-of-memory failures that arise when a uniformly fine grid is applied to large-scale instances.

\paragraph{Tighter initial relaxation.}
By precluding inverse arcs, the initial relaxation model operates over a solution space that is already free of the most egregious chronological violations. This yields a substantially tighter lower bound from the very first iteration, reducing the number of refinement iterations required for convergence.

\subsection{Itinerary Aggregation}
The full model typically exhibits a large optimality gap, as its LP relaxation allows fractional resource flows to be split across multiple adjacent departure-time arcs.
This enables a single resource unit to be reused fractionally, collecting revenue from several nearby time-arcs.
Moreover, some time-arcs generate revenue only when the corresponding resource flow increases from zero to a small positive value.
As a result, the LP relaxation can exploit the high initial revenue portions of multiple time-arcs without paying the full operating cost on any individual arc, leading to significant revenue overestimation.

The relaxation model based on time-arc aggregation we have introduced above, in which each time-arc represents a larger time interval is a valid relaxation of the original model and is designed to be solved efficiently, is not introduced to address the revenue-overestimation issue. The structural revenue-overestimation problem caused by fractional resource reuse is inherited by the relaxation model.
Moreover, due to the nature of the aggregation, this issue may persist even in the integer solutions of the relaxation model.
Consequently, although the relaxation is computationally attractive, it does not by itself eliminate the artificial revenue gains arising from fractional resource flows.

The resulting relaxation value remains close to that of the LP relaxation of the full model, which still permits fractional resource flows to exploit multiple revenue-generating opportunities.
This issue is therefore difficult to resolve by refining the time discretization alone.

\begin{figure}[htbp]
    \centering
    \includegraphics[width=\linewidth]{figure/illustrate_rev_overestimation/rev_distribution}
    \caption{Illustration of Revenue Overestimation and Persistency in Discretization Aggregation}
    \label{fig:rev_overestimation}
\end{figure}

Figure~\ref{fig:rev_overestimation} provides a visual demonstration of this overestimation phenomenon using a specific case. In the diagram, blue arrows represent service assignments, orange bars represent the captured passengers, and gray lines (in Panel 3) indicate the underlying original services.
In the Base Case (Panel 1), the optimal integer solution selects one service $f_0$ ($x=1$), which captures the yield portion of the service (the first third), generating revenue $r_0$.
In the LP Relaxation (Panel 2), the solver exploits fractional variables by assigning three services $f_1, f_2, f_3$ each at value $x=\frac{1}{3}$. Thus each fractional assignment includes the first third of the service, resulting in a revenue of $3r_0$. This erroneously implies that we can serve all potential demand with just one resource unit equivalent of capacity spread over time.
Panel 3 illustrates that our relaxation model defined above inherits this flaw when time points $t_1, t_2, t_3$ are merged into a single window. The model might allow the aggregated service $f_{agg}$ ($x=1$) to claim the summed demand of the underlying original services, again resulting in $3r_0$ revenue as the same way as in the LP relaxation of the full model.

However, under our proposed Itinerary Aggregation, the situation depicted in Panel 3 is eliminated.
The super-itinerary construction explicitly checks the feasibility of combining original itineraries.
Since the departures at $t_1, t_2, t_3$ are too close to be legally operated by a single resource unit (violating minimum separation time), they form an illegal combination for a single super-itinerary.
Our method restricts the attraction summation to valid subsets of original itineraries (using the $\departLimit$ parameter defined in Equation~\ref{def:relax.departLimit}), thereby stripping out the artificial revenue gains and restoring a realistic objective value.

To address this challenge, we introduce an \emph{itinerary aggregation} procedure that complements the time-arc aggregation.

\paragraph{Itinerary aggregation procedure.}
Given a time-arc aggregation defined by the discretization $\discretization$, we further perform an itinerary-level aggregation that is consistent with the aggregated time-arcs.
Let $\itinSet$ denote the original set of itineraries in the full network.
Each itinerary $r \in \itinSet$ consists of an ordered sequence of legs, and each leg uses a time-arc $(o,d,t)$ in the original discretization. Every original itinerary $r$ can thus be represented as a path
\begin{equation}
\origPath{r} = \bigl( (o_1,d_1,t_1), (o_2,d_2,t_2), \dots \bigr)
\end{equation}
where $(o_k,d_k,t_k)$ denotes the origin, destination, and departure time of leg $k$ in itinerary $r$. We refer to this ordered tuple as the \emph{original path signature} of itinerary $r$.

Under the aggregated discretization, each original departure time $t$ is mapped to the beginning of its corresponding time window.
As a result, every original itinerary $r$ induces a discretized path
\begin{equation}
\pathSig{r} = \bigl( (o_1,d_1,\bar t_1), (o_2,d_2,\bar t_2), \dots \bigr)
\end{equation}
where $\bar t_k$ denotes the aggregated departure time of leg $k$.
We define the \emph{discreted path signature} of itinerary $r$ as the ordered tuple $\pathSig{r}$.
Throughout, the notation $(o,d,\bar{t})\in\pathSig{r}$ indicates that the leg with origin $o$, destination $d$, and aggregated departure time $\bar{t}$ occurs as an entry in the tuple; if the same spatial segment appears on multiple legs of the itinerary, each occurrence is counted separately in any summation or product indexed over $\pathSig{r}$.

The itinerary aggregation is then performed by partitioning $\itinSet$ into equivalence classes $\eqClassByPathSig$ based on discreted path signature.
Two itineraries $r,r' \in \itinSet$ are assigned to the same group if and only if they share the same path signature, i.e.,
\[
r \eqRelByPathSig r' \iff \pathSig{r} = \pathSig{r'}
\]
For each equivalence class, we construct a \emph{super-itinerary} that represents all original itineraries in this class.
The super-itinerary inherits the common node sequence and aggregated time-arcs, and replaces all original itineraries in the relaxation model. The set of super-itineraries is denoted by $\itinSetSuper$. A super-itinerary $R$ can be represented by an aggregated path signature in the
same way as an original itinerary. We denote this signature by $\pathSig{R}$.

Attributes associated with itineraries are aggregated conservatively.
Specifically, for each fare class, the fare of a super-itinerary are set to the maximum value among all original itineraries in the corresponding equivalence class.
\begin{equation}
\fare_{r\ell} = \max_{r' \in \eqClassByPathSig_{(r)}} \fare_{r'\ell} \quad \forall r \in \itinSetSuper, \; \ell \in \fareClassSet
\end{equation}

The attraction values of super-itineraries are set to the maximum summation of attractions among all possible combinations of original itineraries in the corresponding equivalence class. This ensures that the super-itinerary's attraction is at least as high as any original itinerary combination it represents. Original itineraries using the same path in the original network can be operated simultaneously without conflict. Therefore, like partitioning $\itinSet$ into equivalence classes $\eqClassByPathSig$ by discreted path signature, for each super-itinerary $R$ we also partition $\eqClassByPathSig_{(R)}$ into equivalence classes $\eqClassByOrigPath^{R}$ by original path signature, two itineraries $r,r' \in \eqClassByPathSig_{(R)}$ are assigned to the same group if and only if they share the same original path signature, i.e.,
\[r \eqRelByOrigPath r' \iff \origPath{r} = \origPath{r'}\]
Given passenger type $p$ and fare class $\ell$. For each equivalence class $\eqClassByOrigPath^{R}_r$, we have a summation of attractions among all original itineraries in the corresponding equivalence class $\sum_{r\in \eqClassByOrigPath^{R}_r} \attr_{rp\ell}$. And now we set the attraction values of a super-itinerary $r$ as the first $\departLimit$ maximum summation among all equivalence class in the quotient set $\eqClassByPathSig_{(R)} /\eqRelByOrigPath$ as follows:
\begin{equation*}
    \begin{aligned}
        \maxAttrCombSet_{Rp\ell}
        & \;=\;
        \arg\max_{\itinCombSet: \itinCombSet \subseteq \eqClassByPathSig_{(R)} /\eqRelByOrigPath,\; |\itinCombSet| \le \combLimit_R}
        \sum_{\mathscr{O} \in \itinCombSet} \sum_{r\in \mathscr{O}} \attr_{rp\ell}
        \\
        \combLimit_R
        & \;=\;
        \prod_{(o,d,\bar{t}) \in \pathSig{R}}
        \min\{\departLimit_{od\bar{t}}, \freq_{o,d}\}
    \end{aligned}
\end{equation*}
\begin{equation}
\attrSuper_{Rp\ell}
\;=\;
\sum_{\mathscr{O} \in \maxAttrCombSet_{Rp\ell}} \sum_{r\in \mathscr{O}} \attr_{rp\ell}
% temp-command \tag{def$_{\attrSuper_{Rp\ell}}$}
\label{def:relax.attrSuper}
\end{equation}
\begin{equation}
\adjAttrSuper_{Rp\ell}
\;=\;
\attrSuper_{Rp\ell} \times \min_{r: r\in\mathscr{O}, \mathscr{O} \in \maxAttrCombSet_{Rp\ell}} \left\{\frac{\adjAttr_{rp\ell}}{\attr_{rp\ell}}\right\}
% temp-command \tag{def$_{\adjAttrSuper_{Rp\ell}}$}
\label{def:relax.shadowAttrSuper}
\end{equation}
where $\departLimit_{od\bar{t}}$ is defined in (\ref{def:relax.departLimit}), $\maxAttrCombSet_{Rp\ell}$ is the set of equivalence classes of original itineraries with the same original path signature that maximizes the total attraction for passenger type $p$ and fare class $\ell$, subject to the constraint that the number of selected equivalence classes does not exceed $\departLimit$. For each leg $(o,d,\bar{t})$ along the super-itinerary's path, the number of services that can physically coexist within the aggregated time window is bounded by $\min\{\departLimit_{od\bar{t}}, \freq_{o,d}\}$, reflecting both the separation-time capacity and the segment frequency. Because original itineraries with different path signatures may consolidate on certain legs, the overall limit $\departLimit$ is taken as the maximum of these per-leg bounds: as long as there exists at least one leg that can accommodate $\departLimit$ distinct services, the corresponding equivalence classes may coexist legally.
This ensures that the relaxation model does not underestimate passenger utility and the objective value while preventing multiple fractional flows from exploiting multiple nearby departure times within the same time window.
\begin{figure}[htbp]
    \centering
    \includegraphics[width=\linewidth]{figure/illustrate_valid_combinations/valid_combinations}
    \caption{Mechanism of Lower Bound Preservation in Itinerary Aggregation}
    \label{fig:valid_combinations}
\end{figure}

Figure~\ref{fig:valid_combinations} illustrates why the itinerary aggregation maintains the lower bound property.
% Consider an equivalence class of original itineraries mapping to the same super-itinerary. These original itineraries (e.g., $r_1, r_2, r_3, r_4, r_5$) may have mutual conflicts due to separation time constraints in the full model. For instance, $r_1$ and $r_3$ cannot be simultaneously operated if they are too close.
% Conceptually (Method 1), the super-itinerary's attraction $\attrSuper$ should be calculated by preventing all conflicts and finding the subset that maximizes the total attraction.
% In the example, the valid subset $\{(r_1, r_2), (r_5)\}$ yields a total attribute of 0.5, which is greater than $\{(r_3, r_4)\}$'s 0.4.
Consider an equivalence class of original itineraries mapping to the same super-itinerary. These original itineraries (e.g., $r_1, r_2, r_3$) may have mutual conflicts due to separation time constraints in the full model. For instance, $r_1$ and $r_2$ cannot be simultaneously operated if they are too close.
Conceptually (Method 1), the super-itinerary's attraction $\attrSuper$ should be calculated by preventing all conflicts and finding the subset that maximizes the total attraction.
In the example, the valid subset $\{(r_1), (r_3)\}$ yields a total attribute of 0.3, which is greater than $\{(r_2)\}$'s 0.2.

However, finding the exact maximum weight independent set considering all the conflicts may be computationally expensive.
Therefore, in our proposed model, we employ a simplified and efficient approximation (Method 2). For each leg $(o,d,\bar{t})$ along the path, the per-leg capacity is bounded by $\min\{\departLimit_{od\bar{t}}, \freq_{od}\}$, accounting for both the separation-time limit (Equation~\ref{def:relax.departLimit}) and the segment frequency. The overall $\departLimit$ is then derived from these per-leg bounds, reflecting the fact that itineraries with different original path signatures may consolidate on shared legs.
We approximate the maximum feasible attraction by selecting the top $\departLimit$ equivalence classes from the quotient set.
Since $\departLimit$ is an upper bound on the size of any valid service sequence, the attraction sum calculated by Method 2 is guaranteed to be greater than or equal to the attraction of any actual conflict-free subset (Method 1), and thus greater than or equal to any feasible integer solution.
Selecting the super-itinerary in the relaxation model effectively grants this "optimistic" revenue.
For any valid solution feasible in the original problem, its revenue cannot exceed this maximum possible revenue calculated for the super-itinerary. Thus, the relaxation objective (Minimize Cost - Revenue) remains less than or equal to the true optimal objective, preserving the validity of the lower bound.

Constraints involving itinerary-level variables are updated to reflect the aggregation. Specifically, constraints (\ref{constr:main.choice.ratio}, \ref{constr:main.choice.limit}, \ref{constr:main.choice.capacity}, \ref{constr:main.choice.covering}) are modified as follows:
\begin{equation}
    \outAttr_{mp}\, \shareVar_{rp\ell} \;\le\; \attrSuper_{rp\ell}\, \outsideShareVar_{mp}
    \quad \forall m\in \marketSet,\, \forall p\in \psgTypeSet,\, \forall r\in \itinSuperInMarket_{m},\, \forall \ell\in \fareClassSet
    % temp-command \tag{2f}
    \label{constr:relax.choice.ratio}
\end{equation}
\begin{equation}
    \left(\frac{\adjOutAttr_{mp}}{\outAttr_{mp}}\right) \outsideShareVar_{mp}
    \;+\;
    \sum_{r\in \itinSuperInMarket_{m}}\sum_{\ell\in \fareClassSet} 
    \left(\frac{\adjAttrSuper_{rp\ell}}{\attrSuper_{rp\ell}}\right) \shareVar_{rp\ell}
    \;\le\; 1
    \quad \forall m\in \marketSet,\; \forall p\in \psgTypeSet
    % temp-command \tag{2g}
    \label{constr:relax.choice.limit}
\end{equation}
\begin{equation}
    \sum_{m\in \marketSet}\sum_{p\in \psgTypeSet}\sum_{\ell\in \fareClassSet}\sum_{r\in \itinSuperUsingArc{m}{a}} 
    \demand_{mp}\, \shareVar_{rp\ell}\;\le\; \sum_{f\in \fleetSet} \fleetCap_{f}\, \fleetAssignVar_{af}
    \quad \forall a\in \flightArc
    % temp-command \tag{2h}
    \label{constr:relax.choice.capacity}
\end{equation}
\begin{equation}
    \sum_{f\in \fleetSet} \fleetAssignVar_{af} \;\ge\; \shareVar_{rp\ell}
    \quad \forall a\in \flightArc,\; \forall m\in \marketSet,\, \forall r\in \itinSuperUsingArc{m}{a},\, \forall p\in \psgTypeSet,\,\forall \ell\in\fareClassSet
    % temp-command \tag{2i}
    \label{constr:relax.choice.covering}
\end{equation}

Through this itinerary aggregation, all fractional resource flows using different original time-arcs within one aggregated time window are forced to compete for the same super-itinerary.
As a consequence, revenue collected on an aggregated time-arc corresponds to at most $\departLimit_{ijt}$ complete resource flows in the original network, effectively eliminating the artificial revenue inflation caused by fractional resource reuse across adjacent departure times.

% \model{\discretization, \itinSetSuper}, the final relaxation model
We denote the final relaxation model incorporating both time-arc aggregation and itinerary aggregation as $\model{\discretization, \itinSetSuper}$. Denote $\model{\discretization} = \model{\discretization, \itinSet}$ when the itinerary set is same as in the full model.

The model $\model{\discretization, \itinSetSuper}$ is written as follows:
\begin{align}
    \min \;\; &
    \sum_{a\in \flightArc}\sum_{f\in \fleetSet} \cost_{af}\, \fleetAssignVar_{af}
    \;-\;
    \sum_{m\in \marketSet}\sum_{p\in \psgTypeSet}\sum_{r\in \itinInMarket_{m}}\sum_{\ell\in \fareClassSet} 
    \demand_{mp}\,\fare_{r\ell}\, \shareVar_{rp\ell}
    % temp-command \tag{2a}
    \label{objfun:relax}
\\
    \text{s.t.} &
    \sum_{a\in \outArcSet{i,t}} \fleetAssignVar_{af} - \sum_{a\in \inArcSet{i,t}} \fleetAssignVar_{af} = 0 
    \quad \forall (i,t)\in \nodeSet,\; \forall f\in \fleetSet
    % temp-command \tag{2b}
    \label{constr:relax.flow}
\\&
    \eqref{constr:relax.fleet_available},
    \eqref{constr:relax.cooldown} \notag
\\&
    \sum_{f\in \fleetSet}\sum_{a\in \segArcSet{s}} \fleetAssignVar_{af} = \freq_{s}
    \quad \forall s\in \segSet
    % temp-command \tag{2e}
    \label{constr:relax.seg_freq}
\\&
    \eqref{constr:relax.choice.ratio},
    \eqref{constr:relax.choice.limit},
    \eqref{constr:relax.choice.capacity},
    \eqref{constr:relax.choice.covering},
    \eqref{constr:relax.domain.x.flight} \notag
\\&
    \fleetAssignVar_{af}\ge 0 \quad \forall a\in \holdingArc, \; \forall f\in \fleetSet
    % temp-command \tag{2k}
    \label{constr:relax.domain.x.holding}
\\&
    \outsideShareVar_{mp}\ge 0 \quad \forall m\in \marketSet, \; \forall p\in \psgTypeSet
    % temp-command \tag{2l}
    \label{constr:relax.domain.w.out}
\\&
    \shareVar_{rp\ell}\ge 0 \quad \forall r\in \itinSet, \; \forall p\in \psgTypeSet, \; \forall \ell\in \fareClassSet
    % temp-command \tag{2m}
    \label{constr:relax.domain.w}
\end{align}

The following theorem guarantees the validity of the relaxation model $\model{\discretization, \itinSetSuper}$ as a lower bound for the original model $\model{\discretizationFull}$.
\begin{theorem}
\label{pro.valid_relaxation}
    松弛模型 $\model{\discretization, \itinSetSuper}$ 是完整模型 $\model{\discretizationFull}$ 的有效松弛。对于 $\model{\discretizationFull}$ 的任意可行解 $\bar{\solution}=(\bar{\fleetAssignVar}, \bar{\shareVar}, \bar{\outsideShareVar})$，存在对应的解 $\bar{\solution}$ 到 $\solution$，且目标值不劣于原解。从 $(\bar{\fleetAssignVar},\bar{\shareVar})$ 到 $(\fleetAssignVar, \shareVar, \outsideShareVar)$ 的映射如下：
    \begin{equation}
    \begin{aligned}
        \fleetAssignVar_{af} &= \sum_{a^\prime\in\Phi_{\arcSet}^{-1}(a)} \bar{\fleetAssignVar}_{a^\prime f},
        \quad \forall f\in \fleetSet,\; \forall a\in \arcSet\\
        \shareVar_{rp\ell} &= \sum_{r' \in \eqClassByPathSig_{(r)}} \bar{\shareVar}_{r'p\ell} \quad \forall r \in \itinSetSuper,\, \forall p \in \psgTypeSet,\, \forall \ell \in \fareClassSet\\
        \outsideShareVar_{mp} &= \bar{\outsideShareVar}_{mp} \quad \forall m \in \marketSet,\, \forall p \in \psgTypeSet
    \end{aligned}
    \label{eq:map_ori2rel}
    \end{equation}
其中 $\eqClassByPathSig_{(r)}$ 是基于路径签名映射到超级行程 $r$ 的原始行程等价类。
\end{theorem}


% % ---------------------------------------- Re model ----------------------------------------
% % plan 1: time potential -> eliminate internal loops
% \input{chapter/lbm/time_potential.tex}
% % plan 2: limit discretization -> ensure no inverse-arc
% 本附录给出聚合时空网络的严格构造规则以及初始离散化的拓扑感知设计。

\paragraph{聚合时空网络的构造规则。}
给定工作离散化$\discretization$，聚合时空网络$\net = (\nodeSet, \arcSet)$定义为$\nodeSet = \{(i,t)\mid \forall i\in\airportSet, t\in\discretization_i\}$，$\arcSet = \holdingArc \cup \flightArc$，满足以下性质：
\begin{property}[离散化$\discretization$的域]
    $\forall i\in\airportSet,\,\{\planningStart\}\subseteq\discretization_i\subseteq\discretizationFull_i$
\end{property}
\begin{property}[节点 $\nodeSet = \Phi_{\nodeSet}(\netFlat, \discretization)$]
    节点集$N$包含节点$(i,\planningStart)\;\forall i\in\airportSet$
\end{property}
\begin{property}[持有弧]
    \label{def:relax.holdingArc}
    每条地面弧$a\in\holdingArc$定义为$a=((i,t),(i,t^{+}))$，其中
    $t^{+}=\begin{cases}
            \min\{\bar{t}\mid\bar{t}>t,\bar{t}\in\discretization_i\} & ,\exists\bar{t}>t,\bar{t}\in\discretization_i\\
            \planningStart & ,\text{otherwise}
    \end{cases}$
\end{property}
\begin{property}[服务弧]
    \label{def:relax.flightArc}
    每条服务弧$a\in\flightArc$定义为$a=((i,t),(j,t^\prime)),\;\forall (i,j)\in\segSet,t\in\discretization_i$，其中
    $t^\prime=\max\{\bar{t}\mid\bar{t}\leq(t+\travelTime_{ij})\;\text{mod}\;(\planningEnd-\planningStart)+\planningStart,\bar{t}\in\discretization_j\}$
\end{property}

\paragraph{弧映射$\Phi$的严格定义。}
对每个节点$i\in\airportSet$，引入投影
\[
    \phi_i:\discretizationFull_i\to\discretization_i,\qquad 
    \phi_i(\tau)=\max\{\bar{t}\in\discretization_i\mid \bar{t}\le \tau\}.
\]
节点映射为$\Phi_{\nodeSet}((i,\tau))=(i,\phi_i(\tau))$。弧映射$\Phi_{\arcSet}:\full{\arcSet}\to\widetilde{\arcSet}$（其中$\widetilde{\arcSet}:=\arcSet\cup\{\bot\}$）定义为：
\begin{equation}
    \Phi_{\arcSet}(((i,\tau)\rightarrow(j,\tau'))) =
    \begin{cases}
        ((i,\phi_i(\tau))\rightarrow(j,\phi_j(\tau'))), & \text{if } i\neq j,\\[2pt]
        ((i,\phi_i(\tau))\rightarrow(i,\phi_i(\tau'))), & \text{if } i=j,\ \phi_i(\tau)<\phi_i(\tau'),\\[2pt]
        \bot, & \text{if } i=j,\ \phi_i(\tau)=\phi_i(\tau').
    \end{cases}
\end{equation}
每条服务弧将其尾部保持在向下取整的聚合时间上，并将其头部设置为旅行时间不增加的最晚聚合时间点；持有弧桥接到下一个聚合节点，若两端点投影到同一聚合节点则折叠为$\bot$。

\paragraph{初始离散化与逆弧消除。}
为给行程聚合提供一个无逆弧的初始聚合网络，本文采用拓扑感知的异构初始离散化。性质~\ref{def:relax.holdingArc}--\ref{def:relax.flightArc}中定义的弧映射规则揭示了均匀粗糙离散化的结构脆弱性。在性质~\ref{def:relax.flightArc}~下，在航段$(j,i)$上从$t \in \discretization_j$出发的服务弧的到达时间被设置为$\discretization_i$中不晚于$t + \travelTime_{ji}$的最晚网格点（下取整映射）。如果对所有节点应用均匀步长$\discretizationStep$且对某些航段$(j,i)$有$\discretizationStep > \travelTime_{ji}$，则映射的到达时间$t'$可能满足$t' \le t$，产生资源单位看起来到达不晚于出发的逆弧。为从源头消除逆弧，每个节点$i \in \airportSet$获得一个定义为最小入站旅行时间的个体初始步长$\initStep_i$：
\begin{equation}\label{eq:init_step}
    \initStep_i = \min_{j:\,(j,i)\in\segSet} \left\{ \travelTime_{ji} \right\}.
\end{equation}
对于没有入站航段的节点，我们设置$\initStep_i = \planningEnd - \planningStart$作为安全默认值。节点$i$的初始时间点集合随后生成为等差数列
\begin{equation}\label{eq:init_discretization}
    \discretization_i = \left\{ t \;\middle|\; t = \planningStart + k \cdot \initStep_i,\quad k \in \mathbb{Z}_{\ge 0},\quad t < \planningEnd \right\}.
\end{equation}

\begin{proposition}[消除逆弧]\label{prop:no_inverse}
    在由\eqref{eq:init_step}--\eqref{eq:init_discretization}定义的初始离散化下，松弛网络$\net$不包含逆弧，即对于每条服务弧$a = ((j,t) \to (i,t')) \in \flightArc$且$(j,i) \in \segSet$，有$t' > t$。
\end{proposition}
\begin{proof}
    由性质~\ref{def:relax.flightArc}，$t' = \max\{\bar{t} \in \discretization_i \mid \bar{t} \le t + \travelTime_{ji}\}$。由于$\discretization_i$是公差为$\initStep_i \le \travelTime_{ji}$（由\eqref{eq:init_step}）的等差数列，区间$(t,\, t + \travelTime_{ji}]$的长度为$\travelTime_{ji} \ge \initStep_i$，因此包含$\discretization_i$的至少一个网格点。故$t' \ge t + \initStep_i > t$（注意$\initStep_i > 0$由构造保证）。
\end{proof}

异构步长自然适应真实世界服务网络的枢纽辐射结构。接收许多短途服务的枢纽节点获得较小的$\initStep_i$，因此获得保持时间精度的更细网格。主要由长途航段服务的辐射节点获得较大的$\initStep_i$和相应更粗的网格。这种自适应行为在不牺牲网络物理有效性的情况下保持初始时间点总数紧凑，避免了当均匀细密网格应用于大规模实例时出现的内存溢出故障。通过排除逆弧，初始松弛模型从第一次迭代就产生了显著更紧的下界，减少了收敛所需的细化迭代次数。

% % plan 3: limit discretization -> ensure no inverse-arc
% \input{chapter/lbm/dynamic_positioning_discovery.tex}
% % ------------------------------------------------------------------------------------------

\section{Upper Bound Generation}
\label{sec:ub}

While the LBM provides a valid lower bound, its formulation on a coarse time-space network often results in solutions that are physically infeasible when mapped back to exact operating times. We propose an MIP-based upper bound (UB) generation module that acts as a repair algorithm based on lower bound information. This approach is designed to reconstruct strict, high-quality feasible solutions from the relaxed state. This section details the UB module within our \algoNotation framework. Unlike purely data-driven heuristics that blindly prune active arcs based on pool statistics, we propose a structure-aware macro-locking strategy called Frequency Fixing, carefully enhanced by rigorous Itinerary Cuts. The upper bound model $\ubm(\discretizationFull,\model{\discretization, \itinSetSuper})$ is formulated on the full time-space network $\full{\net}$ to ensure that the generated solution is directly feasible for the original problem, thus providing a valid upper bound.

\subsection{Frequency Fixing Constraint}
The core assumption of Frequency Fixing is that although the relaxation solution is defined on a coarse network—and thus potentially physically infeasible when mapped to exact timings—it makes highly accurate decisions regarding resource allocation. Therefore, the resource frequency determined by the relaxation model is highly trustworthy and can be locked in the current iteration.

Let $\fleetAssignVar^{*}$ be the integer optimal solution obtained from the relaxation model. We calculate the service frequency for each segment $s=(i,j)$ and resource type $f$:
\begin{equation}
    \freq_{s,f} = \sum_{a \in \segArcSet{s}} \fleetAssignVar_{a,f}^{*}  \quad \in \mathbb{Z}_{\ge 0}
\end{equation}

We then construct the Upper Bound Model by enforcing these counts. We replace the generic frequency constraint~\eqref{constr:main.seg_freq} in the original model $\model{\discretizationFull}$ with:
\begin{equation}
    \sum_{a \in \full{\segArcSet{s}}} \fleetAssignVar_{a,f} = \freq_{s,f}, \quad \forall s \in \segSet, f \in \fleetSet
    % temp-command \tag{3e}
    \label{constr:ubm.freq_fix}
\end{equation}

By forcing the solver to respect the volume of the relaxation solution, the solver's burden is reduced to finding the best feasible timings that satisfy flow balance. This effectively reduces the multi-faceted original model into a pure scheduling problem.

\subsection{Addressing Revenue Overestimation via Itinerary Cuts}
The standard Frequency Fixing only locks macro-level resource frequencies but leaves the revenue overestimation problem, which makes the model difficult to solve. To mitigate this, we inject super itineraries from the relaxation model into the upper bound model. We introduce continuous market share variables $\shareSVar_{R,p,\ell} \in [0, 1]$ in the UB model to represent the proportion of the market demand choosing super itinerary $R$ for passenger type $p$ and fare class $\ell$. By adding cuts that link the super itinerary share variables $\shareSVar_{R,p,\ell}$ with the original fine-grained share variables $\shareVar_{r,p,\ell}$, we enforce that the revenue contribution of a super itinerary must dominate the total revenue of its constituent fine-grained itineraries, thereby significantly tightening the upper bound and accelerating convergence of the branch-and-bound search. The detailed formulation of these cuts is provided below:
\begin{equation}
    \outAttr_{mp}\, \shareSVar_{Rp\ell} \;\le\; \attrSuper_{Rp\ell}\, \outsideShareVar_{mp}
    \quad \forall m\in \marketSet,\, \forall p\in \psgTypeSet,\, \forall R\in \itinSuperInMarket_{m},\, \forall \ell\in \fareClassSet
    % temp-command \tag{3f}
    \label{constr:ubm.itinerary_cut.choice.ratio}
\end{equation}
\begin{equation}
    \left(\frac{\adjOutAttr_{mp}}{\outAttr_{mp}}\right) \outsideShareVar_{mp}
    \;+\;
    \sum_{R\in \itinSuperInMarket_{m}}\sum_{\ell\in \fareClassSet} 
    \left(\frac{\adjAttrSuper_{Rp\ell}}{\attrSuper_{Rp\ell}}\right) \shareSVar_{Rp\ell}
    \;\le\; 1
    \quad \forall m\in \marketSet,\; \forall p\in \psgTypeSet
    % temp-command \tag{3g}
    \label{constr:ubm.itinerary_cut.choice.limit}
\end{equation}
\begin{equation}
    \sum_{m\in \marketSet}\sum_{p\in \psgTypeSet}\sum_{\ell\in \fareClassSet}\sum_{R\in \itinSuperUsingArc{m}{a}} 
    \demand_{mp}\, \shareSVar_{Rp\ell}\;\le\; \sum_{a'\in \Phi^{-1}(a)}\sum_{f\in \fleetSet} \fleetCap_{f}\, \fleetAssignVar_{a'f}
    \quad \forall a\in \flightArc
    % temp-command \tag{3h}
    \label{constr:ubm.itinerary_cut.choice.capacity}
\end{equation}
\begin{equation}
    \sum_{a'\in \Phi^{-1}(a)}\sum_{f\in \fleetSet} \fleetAssignVar_{a'f} \;\ge\; \shareSVar_{Rp\ell}
    \quad \forall a\in \flightArc,\; \forall m\in \marketSet,\, \forall R\in \itinSuperUsingArc{m}{a},\, \forall p\in \psgTypeSet,\,\forall \ell\in\fareClassSet
    % temp-command \tag{3i}
    \label{constr:ubm.itinerary_cut.choice.covering}
\end{equation}
\begin{equation}
   \shareSVar_{R,p,\ell} \cdot \fare_{R,\ell} \;\le\; \sum_{r \in R} \shareVar_{r, p, \ell} \cdot \fare_{r,\ell}
   % temp-command \tag{3n}
   \label{constr:ubm.itinerary_cut.cut}
\end{equation}

Collectively, constraints \eqref{constr:ubm.itinerary_cut.choice.ratio}-\eqref{constr:ubm.itinerary_cut.choice.covering} make sure that the super itinerary share $\shareSVar_{R,p,\ell}$ represents the customer choice behavior at the macro level (time windows in the partial network) and yields no less share than the aggregated share of its constituent fine-grained itineraries. The cut \eqref{constr:ubm.itinerary_cut.cut} explicitly links the super itinerary share variable with the fine-grained share variables, ensuring that the revenue contribution of the super itinerary dominates the sum of its underlying itineraries. By enforcing these relationships between the super itinerary and its constituent itineraries, we effectively tighten the upper bound and guide the solver towards more promising regions of the solution space, thus accelerating convergence to the global optimum.

Let $\textbf{c}^{\mathrm{LB}}$ denote the cost term of optimal objective value of $\model{\discretization, \itinSetSuper}$. Let $\textbf{r}^{\mathrm{LB}}$ denote the revenue term of optimal objective value of $\model{\discretization, \itinSetSuper}$. By Proposition~\ref{pro.valid_relaxation}, $\model{\discretization, \itinSetSuper}$ constitutes a valid relaxation of the full model; it follows that the optimal revenue of any feasible solution of the original problem is bounded below by $\textbf{r}^{\mathrm{LB}}$. We exploit this bound explicitly by appending the following valid inequality:
\begin{equation}
    \sum_{m\in \marketSet}\sum_{p\in \psgTypeSet}\sum_{\ell\in \fareClassSet}\sum_{r\in \itinInMarket_m}
    \demand_{mp}\, \fare_{r\ell}\, \shareVar_{rp\ell} \;\le\; \textbf{r}^{\mathrm{LB}}
    % temp-command \tag{3o}
    \label{constr:ubm.obj_cut}
\end{equation}
This cut explicitly excludes all LP solutions whose total revenue falls strictly below the known lower bound, thereby tightening the feasible region and reducing the effective search space of the branch-and-bound procedure.

Thus, we have a upper bound generation module that synergistically combines frequency fixing and itinerary cuts, effectively repairing the relaxed LB solution into a high-quality feasible solution that guides the exact solver towards global optimality, which is written as follows:
\begin{equation}\begin{aligned}
    \ubm(\discretizationFull,\model{\discretization, \itinSetSuper}) = \min_{\shareSVar, \shareVar, \fleetAssignVar} \quad & \textbf{c}^{\mathrm{LB}} - \sum_{m\in \marketSet}\sum_{p\in \psgTypeSet}\sum_{\ell\in \fareClassSet}\sum_{r\in \itinInMarket_m} \demand_{mp}\, \shareVar_{rp\ell} \fare_{r,\ell}
    % temp-command \tag{3a}
\\
    \text{s.t.} \quad
    &
    \eqref{constr:main.flow}, \eqref{constr:main.fleet_available}, \eqref{constr:main.cooldown}, \eqref{constr:main.seg_freq} \notag
\\&
    \eqref{constr:main.choice.ratio}, \eqref{constr:main.choice.limit}, \eqref{constr:main.choice.capacity}, \eqref{constr:main.choice.covering} \notag
\\&
    \eqref{constr:main.domain.x.flight}, \eqref{constr:main.domain.x.holding}, \eqref{constr:main.domain.w.out}, \eqref{constr:main.domain.w} \notag
\\&
    \eqref{constr:ubm.freq_fix}, \eqref{constr:ubm.itinerary_cut.choice.ratio}, \eqref{constr:ubm.itinerary_cut.choice.limit}, \eqref{constr:ubm.itinerary_cut.choice.capacity}, \eqref{constr:ubm.itinerary_cut.choice.covering}, \eqref{constr:ubm.itinerary_cut.cut}, \eqref{constr:ubm.obj_cut} \notag
\end{aligned}\end{equation}
The following theorem guarantees the validity of the upper bound generated by $\ubm(\discretizationFull,\model{\discretization, \itinSetSuper})$ for the original model $\model{\discretizationFull}$.
\begin{theorem}
    \label{thm:valid_upper_bound}
    对于任意 $\discretization, \itinSetSuper$，若模型 $\ubm(\discretizationFull, \model{\discretization, \itinSetSuper})$ 具有非空可行域，则其最优目标值是原始模型 $\model{\discretizationFull}$ 的上界。
\end{theorem}
% --- Explicit note about UBM infeasibility handling ---
% If locked aggregated frequencies cannot be realized on the full time-space network,
% UBM is infeasible; in our algorithm this corresponds to an objective of +\infty
% and the candidate upper bound is ignored by the master algorithm (Algorithm~1).
% We state this explicitly to clarify that no additional fallback repair is required:
If the UB model $\ubm(\discretizationFull,\model{\discretization, \itinSetSuper})$ becomes infeasible under the locked aggregated frequencies $\freq_{s,f}$, the solver effectively returns an outcome (treated as an upper bound of $+\infty$ in our algorithm logic). In such cases the algorithm does not update the incumbent UB and proceeds to the refinement phase driven by LBM-infeasibilities.
\subsection{Variable Reduction}
To accelerate the solution process, we apply a series of exact variable reductions based on the fixed aggregated frequencies $\freq_{s,f}$. Because $\freq_{s,f}$ dictates the available network capacity, any segment with zero assigned services fundamentally breaks connectivity, allowing us to safely fix many variables to zero and eliminate redundant downstream variables and constraints prior to sending the model to the solver.

For any segment $s \in \segSet$ and resource type $f \in \fleetSet$, if the locked frequency $\freq_{s,f} = 0$, it implies that no services of type $f$ can be operated on segment $s$. Consequently, we immediately cascade this zero assignment to all underlying service assignment variables:
\begin{equation}
    \fleetAssignVar_{a,f} = 0, \quad \forall a \in \full{\segArcSet{s}} \text{ if } \freq_{s,f} = 0
\end{equation}
This satisfies the frequency constraint \eqref{constr:ubm.freq_fix} intrinsically for these combinations, allowing us to remove these variables and constraints from the domain.

Since an itinerary $r$ (or super itinerary $R$) fundamentally consists of a sequence of time-space service arcs, removing all capacity from any of its underlying spatial segments directly severs these physical paths. For notational convenience, we write $s \in r$ (or $s \in R$) to indicate that the spatial segment $s$ is traversed by the time-space arcs of itinerary $r$ (or super itinerary $R$). If an itinerary relies on a segment $s$ where no services of any resource type are assigned (i.e., $\sum_{f\in\fleetSet} \freq_{s,f} = 0$), the itinerary becomes physically impossible. Due to the covering constraints \eqref{constr:main.choice.covering} and \eqref{constr:ubm.itinerary_cut.choice.covering}, the market share associated with such itineraries falls to zero:
\begin{equation}
    \shareVar_{r,p,\ell} = 0, \quad \text{if } \exists s \in r \text{ s.t. } \sum_{f\in\fleetSet} \freq_{s,f} = 0
\end{equation}
\begin{equation}
    \shareSVar_{R,p,\ell} = 0, \quad \text{if } \exists s \in R \text{ s.t. } \sum_{f\in\fleetSet} \freq_{s,f} = 0
\end{equation}
These zero-valued variables can be explicitly removed from the model. As a consequence, the itinerary cut constraint \eqref{constr:ubm.itinerary_cut.cut} for a given super itinerary $R$ can be completely deleted only if all of its constituent fine-grained itineraries $r \in R$ are closed (i.e., fixed to zero). If only a subset of $r \in R$ are removed, the constraint \eqref{constr:ubm.itinerary_cut.cut} must be retained to bind the super itinerary variable $sw_R$ with the remaining active fine-grained itineraries.

For any market $m \in \marketSet$, if all its constituent active itineraries $r \in \itinInMarket_{m}$ and super itineraries $R \in \itinSuperInMarket_{m}$ are eliminated by the reduction above, a purely unserved scenario is triggered. Passengers in this market simply have no available service options to choose from.
According to the choice limit identity constraint \eqref{constr:main.choice.limit}, the outside option share variable $\outsideShareVar_{mp}$ inevitably becomes a deterministic constant bounded entirely by the respective utilities:
\begin{equation}
    \outsideShareVar_{mp} = \frac{\outAttr_{mp}}{\adjOutAttr_{mp}}, \quad \text{if } \shareVar_{r,p,\ell} = 0 \; \forall r \in \itinInMarket_m
\end{equation}
Therefore, we can completely eliminate $\outsideShareVar_{mp}$ as a variable and remove market $m$'s related choice limit \eqref{constr:main.choice.limit} and ratio limits \eqref{constr:main.choice.ratio} entirely from the optimization framework, significantly compressing the constraint matrix dimensions.


\section{Refining Discretization}
\label{sec:refinement}

Because the relaxation model relies on an aggregated time discretization $\discretization$, its optimal solutions may employ services that violate strict real-world constraints. In our algorithm sequence, the Refinement step is invoked whenever such an illegal lower-bound solution is produced. By systematically analyzing network flow, we identify specific chronologically invalid connections (which we denote as ``Impossible Departures''), services that too close to be physically realized (which denote as ``Inter Departures'') and super itineraries whose revenue contribution connot be achieved by any original itinerary combination (which we denote as ``Phantom Revenue''). We then strategically insert new time points to the discretization $\discretization$ to eliminate these infeasibilities. In this section, we first define the violation metrics for each of these three types of violations, and show that existence of any of these violations is necessary and sufficient for the relaxation solution to be illegal in the original model. Then, we present the refinement strategy to eliminate these violations, and prove that the algorithm will converge to a globally optimal solution after finite iterations. Note that the relaxation solutions can be illegal in the original model even when its objective value equals that produced by the upper bound, which means the algorithm is converged. We do not need to further refine the discretization to eliminate the violations in this case because we have already obtained a globally optimal solution.

\begin{definition}
\label{def:illegal_solution}
Let $\solution_{rel} = (\fleetAssignVar^*, \shareVar^*, \outsideShareVar^*)$ be an optimal solution to the relaxation model $\model{\discretization, \itinSetSuper}$ with an objective value of $Z_{rel}^*$. 

The relaxation solution $\solution_{rel}$ is defined as illegal with respect to the full model $\model{\discretizationFull}$ if there does not exist a passenger flow allocation $(\shareVar, \outsideShareVar)$ such that the combined assignment $\solution = (\fleetAssignVar^*, \shareVar, \outsideShareVar)$ is feasible in $\model{\discretizationFull}$ and achieves an objective value $Z = Z_{rel}^*$.
\end{definition}

\subsection{Impossible Departures}
The relaxation model may utilize resource units that physically have not arrived yet but are available in the model for the reason that we allow some short arcs in the relaxation model. This leads to infeasible service connections when we try to implement the solution in reality.
To detect unrealistic service connections, we apply a flow balance check at each active node $(i, t)$ in the time-space network.

At a specific node $(i, t)$, incoming flow from upstream must support outgoing flow. We categorize the incoming arcs of a resource type $f$ into components. Denote $\holdingArcIn(i, t, f)$ as the holding flow into node $(i, t)$ for resource type $f$, $\flightArc[-,legal](i, t, f)$ as services with feasible turn-around times, i.e., arcs $a = (i, t', j, t)$ where $t' + \travelTime_a \le t$, and $\flightArc[-,illegal](i, t, f)$ as short arcs where $t' + \travelTime_a > t$. We have
$$\inArcSet{i,t}=\holdingArcIn(i, t, f)\uplus \flightArc[-,legal](i, t, f) \uplus \flightArc[-,illegal](i, t, f)$$

The outgoing flows from node $(i, t)$ for resource type $f$ consist of service departures and holding flows. Denote $\departArc(i, t, f)$ as the departure flow from node $(i, t)$ for resource type $f$, and $\holdingArcOut(i, t, f)$ as the holding flow out of node $(i, t)$ for resource type $f$. We have
$$\outArcSet{i,t}=\departArc(i, t, f)\uplus \holdingArcOut(i, t, f)$$

At any node $(i, t)$, the available resource units for departure for resource type $f$ is the sum of holding flow and legal service arrivals, i.e., $\sum_{a\in\holdingArcIn(i, t, f)}\fleetAssignVar_{a,f} + \sum_{a\in\flightArc[-,legal](i, t, f)}\fleetAssignVar_{a,f}$. The actual departure demand is $\sum_{a\in\departArc(i, t, f)}\fleetAssignVar_{a,f}$. If the departure demand exceeds the available resource units, i.e.,
\begin{equation}
    \sum_{a\in\departArc(i, t, f)}\fleetAssignVar_{a,f} > \sum_{a\in\holdingArcIn(i, t, f)}\fleetAssignVar_{a,f} + \sum_{a\in\flightArc[-,legal](i, t, f)}\fleetAssignVar_{a,f}
\end{equation}
then we have impossible departures at node $(i, t)$ for resource type $f$. This means the model is relying on illegal service arcs to satisfy the departure demand, which is infeasible in reality.
For each active node $(i, t)$ and resource type $f$, we calculate the illegal departure quantity $\impDeparts(i,t,f)$:
\begin{equation}
    \impDeparts(i,t,f) = \sum_{a\in\departArc(i, t, f)}\fleetAssignVar_{a,f} - \sum_{a\in\holdingArcIn(i, t, f)}\fleetAssignVar_{a,f} - \sum_{a\in\flightArc[-,legal](i, t, f)}\fleetAssignVar_{a,f}
\end{equation}

Here, $\impDeparts(i,t,f)$ represents the net demand for resource units that cannot be satisfied by legal means. If $\impDeparts(i,t,f) > 0$, the model is forced to use illegal arrivals $\flightArc[-,illegal]$ to satisfy this departure demand. This metric directly quantifies the illegality.

By adding time points to the discretization $\discretization$ we can ensure that all these illegal departures will not occur in the refined model. Denote the refined discretization as $\discretization^{1}$. The refine procedure will be described in the Subsection~\ref{sec:refine_process}.
\subsection{Inter-Departure Violations}
In the relaxation model, multiple departures on the same segment may be scheduled within a segment's minimum inter-departure time, which violates operational separation constraints. To detect and quantify such violations we define the departure flow inside the window starting at time $t$ on segment $(i,j)$ and compare it with the allowable limit.

Let the set of arcs on segment $s = (i,j)$ be denoted by $\segArcSet{s}$. For each arc $a\in\segArcSet{s}$ write $a=(i,t_a,j,t'_a)$ so that $t_a$ denotes the departure time of $a$. Using the resource assignment variables $\fleetAssignVar_{a,f}$, we define the aggregated departure flow in the interval $[t, t+\minInterdepTime_{s})$ as
\begin{equation}
    \arcDepartFlow_{ijt} = \sum_{a\in\segArcSet{s}} \mathbf{1}\{t_a \in [t,\, t+\minInterdepTime_{s})\} \sum_{f\in\fleetSet} \fleetAssignVar_{a,f}.
\end{equation}

Here $\arcDepartFlow_{ijt}$ denotes the total number of departures measured in assigned flow on segment $(i,j)$ inside the minimum-separation window starting at $t$. In the full-model scope the instantaneous capacity permits at most one departure per separation window, hence the allowable upper bound is $1$. A separation violation occurs when the actual flow exceeds $1$.

We quantify the excess by the Separation Violation Flow
\begin{equation}
    \sepViolFlow_{ijt} = \max\bigl\{0,\; \arcDepartFlow_{ijt} - 1\bigr\}.
\end{equation}

Equivalently one may define the Separation Violation Amount
\begin{equation}
    \sepViolAmount_{ij}(t) = \arcDepartFlow_{ijt} - 1,
\end{equation}
and truncate to non-negative values when reporting violations.

Detection procedure: compute $\sepViolAmount_{ijt}$ for all segments $(i,j)$ and all discretized start times $t\in\discretization$. If there exists any $t$ such that
\begin{equation}
    \sepViolAmount_{ijt} > 0,
\end{equation}
then the solution contains an inter-departure violation (the second illegal condition in Theorem~\ref{thm:relax_illegal}). The metric identifies which segments and time windows require refinement or constraint tightening. By adding time points to the discretization $\discretization^{1}$ we can ensure that all these inter-departure violations will not occur in the refined model. Denote the refined discretization as $\discretization^{2}$. The refine procedure will be described in the Subsection~\ref{sec:refine_process}.
\subsection{Phantom Revenue}
When the relaxation solution $\solution = (\fleetAssignVar^*, \shareVar^*, \outsideShareVar^*)$ satisfies both resource availability and inter-departure constraints, it might still be illegal if it relies on "Phantom Revenue." This occurs when the revenue claimed by the aggregated super itineraries cannot be physically realized by any valid combination of their constituent original itineraries in the full discretization network.

Due to the highly coupled nature of the Sales-Based Linear Program (SBLP), the outside option share $\outsideShareVar$ acts as a dynamic variable. A local decrease in the attractiveness of available itineraries can shift demand to the outside option, thereby implicitly raising the choice bounds for other itineraries and reallocating passenger flows. Because of this endogeneity, evaluating Phantom Revenue via static slack analysis is insufficient. Instead, we formulate a Revenue Feasibility Subproblem (RFS) to establish a strict necessary and sufficient condition.

Given the fixed physical service assignment $\fleetAssignVar^*$, we seek to detect whether the target revenue contribution of each super itinerary $R$, defined as $Rev_{Rpl}^* = \demand_{mp} \fare_{Rl} \shareVar_{Rpl}^*$, can be attained by assigning valid shares $\shareVar_{rpl}$ to its constituent original itineraries $r \in \eqClassByPathSig_{(R)}$. We introduce continuous slack variables $s_{Rpl} \ge 0$ to capture the revenue shortfall for each super itinerary. The RFS is formulated as the following linear program:
\full{\begin{align}
RFS(\solution)=\min_{\shareVar, \outsideShareVar, s} \quad & \sum_{m\in \marketSet} \sum_{p\in \psgTypeSet} \sum_{l\in \fareClassSet} \sum_{R \in \itinSuperInMarket_m} s_{Rpl} \label{eq:rfs_obj}\\
\text{s.t.} \quad & \sum_{r \in \eqClassByPathSig_{(R)}} \demand_{mp} \fare_{rl} \shareVar_{rpl} + s_{Rpl} = \demand_{mp} \fare_{Rl} \shareVar_{Rpl}^* \quad \forall R \in \itinSuperInMarket_m, p\in \psgTypeSet, l\in \fareClassSet \label{eq:rfs_target}\\
& \outAttr_{mp} \shareVar_{rpl} \le \attr_{rpl} \outsideShareVar_{mp} \quad \forall m\in \marketSet, p\in \psgTypeSet, l\in \fareClassSet, r \in \itinInMarket_m \label{eq:rfs_ratio}\\
& \left(\frac{\adjOutAttr_{mp}}{\outAttr_{mp}}\right) \outsideShareVar_{mp} + \sum_{r \in \itinInMarket_m} \sum_{l\in \fareClassSet} \left(\frac{\adjAttr_{rpl}}{\attr_{rpl}}\right) \shareVar_{rpl} \le 1 \quad \forall m\in \marketSet, p\in \psgTypeSet \label{eq:rfs_limit}\\
& \sum_{m\in \marketSet}\sum_{p\in \psgTypeSet}\sum_{l\in \fareClassSet} \sum_{r \in \itinUsingArc{m}{a}} \demand_{mp} \shareVar_{rpl} \le \sum_{f\in \fleetSet} \fleetCap_f \fleetAssignVar_{af}^* \quad \forall a \in \flightArc \label{eq:rfs_cap}\\
& \shareVar_{rpl} \ge 0, \quad \outsideShareVar_{mp} \ge 0, \quad s_{Rpl} \ge 0 \nonumber
\end{align}}

Constraint (\ref{eq:rfs_target}) forces the underlying original itineraries to meet the relaxation revenue, using $s_{Rpl}$ to absorb any unavoidable deficit. Constraints (\ref{eq:rfs_ratio})-(\ref{eq:rfs_cap}) ensure that the recovered passenger flow strictly adheres to the physical capacity of $\fleetAssignVar^*$ and the original discrete choice bounds. 

The relaxation solution $\solution$ contains Phantom Revenue (and is therefore illegal) \textit{if and only if} the optimal objective value of the RFS is strictly greater than zero (i.e., $\sum s_{Rpl}^* > 0$). The time slots (i.e. a set of segment-time pairs) corresponding to the super itineraries with positive slack $s_{Rpl}^* > 0$ are the culprits that contribute to the Phantom Revenue. By adding time points to the discretization $\discretization^{2}$ we can ensure that all these phantom revenue violations will not occur in the refined model. Denote the refined discretization as $\discretization^{3}$. The refine procedure will be described in the Subsection~\ref{sec:refine_process}.
% ========================

The following theorem guarantees that the relaxation solution is illegal in the original model if and only if there are at least one of these three violations.

\begin{theorem}
\label{thm:relax_illegal}
A solution $\solution = (\fleetAssignVar^*, \shareVar^*, \outsideShareVar^*)$ to the relaxation model $\model{\discretization}$ is illegal in the original model $\model{\discretizationFull}$ if and only if one of the following three conditions holds:
\begin{enumerate}
    % \textbf{Impossible Departure}: 
    \item There exists at least one node $(i, t)$ and resource type $f$ such that the departure demand at $(i, t)$ for resource type $f$ exceeds the available supply, i.e., $\impDeparts(i,t,f) > 0$.
    % \textbf{Inter Departure}: 
    \item There exist a segment $(i,j)$ and departure time $t$ such that more than one service departs from node $i$ to $j$ within the minimum turn-around time, i.e., $\arcDepartFlow_{ijt} > 1$.
    % \textbf{Phantom Revenue}: 
    \item There exists a super itinerary in the solution whose revenue contribution cannot be achieved by any combination of original itineraries, i.e., $RFS(\solution) > 0$.
\end{enumerate}
\end{theorem}

% ========================

\subsection{Refinement Strategy}\label{sec:refine_process}
\paragraph{Impossible Departure}
If $\impDeparts(i, t, f) := \impDeparts > 0$, we have a violation. The culprits are the set of illegal arrival arcs $\mathcal{K} = \flightArc[-,illegal](i, t, f)$.

For each $k \in \mathcal{K}$, we determine its correction time $t'_{arr} = t_{\text{dep}}(k) + \travelTime_k$, which is the true physical arrival time.  We resolve the violation using the following greedy process. First, sort the culprits $\mathcal{K}$ by their flow volume $\fleetAssignVar_{k}$ in descending order. Then, greedily select culprits from $\mathcal{K}$ until the sum of their flow covers the deficit $\impDeparts$.
Finally, for each selected culprit $k$, add its true arrival time $t'_{arr}$ to the time discretization of node $i$ as $\discretization_i \gets \discretization_i \cup \{t'_{arr}(k)\}$. This ensures that in the next iteration, the model can no longer rely on the illegal arc $k$ to satisfy the departure demand at $(i, t)$, as it will now have a more accurate representation of the resource unit's availability time.

As a result, the time node grid is densified. In the next iteration, the previously problematic arc $k$ will be mapped to land at $t'_{arr}$ (which is later than $t$). This accurately delays the supply of the resource unit, rendering it unavailable for the impossible departure at $t$, thus the current relaxation solution will be infeasible in the next iteration.

\paragraph{Inter Departure}
If an inter-departure violation is detected on segment $(i,j)$ with $\sepViolAmount_{ij}(t) > 0$, it indicates that multiple physical services assigned to this segment are aggregated into an overlapping separation window. To resolve this, we identify the set of active original service arcs $k \in \segArcSet{s}$ whose physical departure times $t_a(k)$ fall within the violating interval $[t, t + \minInterdepTime_{s})$. 
Following a similar greedy logic to the impossible departures, we sort these culprit arcs by their flow volume $\fleetAssignVar_{k}$ and select the top arcs. For each selected arc $k$, we insert its true physical departure time $t_a(k)$ into the discretization of the origin node $i$ as $\discretization_i \gets \discretization_i \cup \{t_a(k)\}$. This densification ensures that the model can distinguish between these close departures in the next iteration, eventually forcing them to satisfy the minimum separation constraint.

\paragraph{Phantom Revenue}
When a super itinerary $R$ is identified as having phantom revenue (i.e., $s_{Rpl}^* > 0$ in the RFS), the current discretization is too coarse to represent the true physical connectivity or capacity constraints of its constituent itineraries $S_{(R)}$. To eliminate this infeasibility, we refine the discretization along the path of $R$. For each original itinerary $r \in S_{(R)}$ that contributes to the revenue shortfall, we extract the physical arrival and departure times of all its constituent service legs. Let $t^{phys}$ be such a physical time point at node $i$. We update the discretization $\discretization_i \gets \discretization_i \cup \{t^{phys}\}$. By splitting the time nodes at the transition points of the super itinerary, we break the equivalence class $S_{(R)}$, forcing the relaxation model to acknowledge the chronological or capacity conflicts that originally led to the revenue overestimation.

The following theorem guarantees that the refinement strategy will converge to a globally optimal solution after finite iterations.
\begin{theorem}
\label{thm:refine_converge}
第~\ref{sec:refine_process}小节中描述的细化策略，在松弛解于原始模型中非法时，将至少向时间离散化 $\discretization$ 中添加1个新时间点。因此，经过有限次迭代后，算法将收敛到全局最优解。
\end{theorem}